% Appendix B (group D): reproduction. Every path below was checked to exist in the repository on 25 November 2026 (H213--H236 folder on 28 November 2026).
\section{Reproduction}
\label{sec:repro}

\paragraph{Source version.} The article reports the Quantum Sphaera Companion \cite{QSC}. The baseline is v5.6.0
(28~November 2026, Version DOI \href{https://doi.org/10.5281/zenodo.24632437}{10.5281/zenodo.24632437}, Concept DOI
\href{https://doi.org/10.5281/zenodo.19210728}{10.5281/zenodo.19210728}; predecessor v5.5.0,
\href{https://doi.org/10.5281/zenodo.24563811}{10.5281/zenodo.24563811}); it contains the research rounds
up to v4.10.0-r75, among them r60, r61 (formula corrections), r62--r65 (the tests after r60 that are cited here: H105,
H128--H156 and H162), r66, r67 (H164, H165 and the counting postulate), r69 (names of the frames and lepton sectors), r70 (frames on
the state sphere), r72 (photon-clock guard restated), r73, r74 (lepton phase, H213--H236) and r75 (dated register of predictions).
The German preprint 0.5.0 \cite{QS-preprint-de} is built on the same Companion version. The sources are in the public repository
\url{https://github.com/Moriarty1982/QuantumSphere-Unfold}. Nothing after r75 is used. In that repository:
\begin{itemize}
\item \nolinkurl{sources/research/v480-companion.tex} is the entry point of the Companion (historical file name);
\item \nolinkurl{sources/research/parts/v490-post-release-layers.tex} holds the rounds v4.9.0-r01 to r86, cited here as
  \textsc{c}\,4.9-rNN;
\item \nolinkurl{sources/research/parts/v4100-post-release-layers.tex} holds the rounds v4.10.0-r01 to r75, cited as
  \textsc{c}\,rNN, with labels \texttt{sec:v4100-rNN-\dots}.
\end{itemize}
Each round ends with a ledger and guard update and names its notes and scripts, in a ``Proof and measurement
sources'' paragraph or, in most later rounds, in the results or in a closing notes paragraph.

\paragraph{Where the scripts live.} The notes of a round are in \nolinkurl{notes/research/v4_10_0_rNN/}. The folder name
follows the working sequence, not the round label: \nolinkurl{v4_10_0_w1/} holds the notes of v4.9.0-r81,
\nolinkurl{v4_10_0_r82/} to \nolinkurl{v4_10_0_r86/} those of v4.9.0-r82 to r86, and \nolinkurl{v4_10_0_r01/} onwards those
of v4.10.0-r01 onwards. Scripts and their recorded outputs are in the subfolder \nolinkurl{llm/} or, for later rounds, in a
topic folder that the round's \texttt{README.md} names (for example \nolinkurl{notes/research/op5_attractor_2026_11_23/} for
H124 and \nolinkurl{notes/research/h125_relatives_cp_2026_11_24/} for H125). Pre-registration and result notes
(\nolinkurl{Vorab_Notiz_H*.md}, \nolinkurl{Ergebnis_H*.md}) are in \nolinkurl{notes/research/predictions/} for H42, H46,
H48--H112 and H124--H137, in the \nolinkurl{llm/} folder of their round for earlier tests, in \nolinkurl{llm/} at the top
level for H113--H123 and in \nolinkurl{notes/research/hilbert_shift_felder_2026_11_24/} for H138--H156 and, with one subfolder per test, for H213--H236
(scripts k426--k451 with Wolfram cross-checks; folder list in \rnd{r73} and \rnd{r74}). The commit of a pre-registration note can be compared with the date of the computation. Table~\ref{tab:repro-act} lists the scripts
behind the numbers of Section~\ref{sec:act}.

\begin{table}[ht]
\centering
\footnotesize
\begin{tabularx}{\textwidth}{@{}>{\raggedright\arraybackslash}p{0.33\textwidth}>{\raggedright\arraybackslash}X l@{}}
\toprule
result & script (path under \nolinkurl{notes/research/}) & tool\\
\midrule
$\mathrm{Im}\,S$ identity of the complex reading & \nolinkurl{v4_10_0_w1/llm/k170_imS_identitaet.py} & Python\\
vacuum decay of an independent T field (H76) & \nolinkurl{v4_10_0_r83/llm/k191_vakuumzerfall.py} & Python\\
second reading on FRW (H80, H81) & \nolinkurl{v4_10_0_r84/llm/k196_zweite_lesart_frw.py} & Python\\
masses under the second reading (H82) & \nolinkurl{v4_10_0_r85/llm/k197_massen_zweite_lesart.py} & Python\\
CPT assignment at the Big Bang (H86) & \nolinkurl{v4_10_0_r85/llm/k201_cpt_zuordnung.py} & Python\\
admixture, second system (H87) & \nolinkurl{v4_10_0_r86/llm/k202_wolfram_gegenprobe.wl} & Wolfram\\
field equations of W1 (H88) & \nolinkurl{v4_10_0_r86/llm/k203_w1_feldgleichungen.wl} & Wolfram\\
energy--momentum tensor of W1 (H89) & \nolinkurl{v4_10_0_r86/llm/k204_w1_einstein.wl} & Wolfram\\
non-linear Einstein--Cartan equation, sign (H90) & \nolinkurl{v4_10_0_r01/llm/k205_w1_einstein_nichtlinear.py},
  output \nolinkurl{k205_out.json} & Python (jets)\\
weak-field sign check & \nolinkurl{v4_10_0_r01/llm/k205b_schwachfeld.wl} & Wolfram\\
fermions at the Big Bang & \nolinkurl{v4_10_0_r03/llm/k209_fermionen_urknall.py}, output \nolinkurl{k209_out.txt} & Python (sympy)\\
the two senses of rotation & \nolinkurl{v4_10_0_r03/llm/k211_symmetrie_drehsinne.py}, output \nolinkurl{k211_out.txt} & Python\\
branch gas on FRW (H95) & \nolinkurl{v4_10_0_r04/llm/k212_zweiggas_frw.py}, output \nolinkurl{k212_out.json} & Python\\
\bottomrule
\end{tabularx}
\caption{Scripts behind Section~\ref{sec:act}. The Wolfram files of \texttt{v4\_10\_0\_r86/\allowbreak llm/\allowbreak } have no separate output
files; their results and the disclosed errors of the checking code are recorded in
\texttt{predictions/\allowbreak Ergebnis\_H87\_*}, \texttt{Ergebnis\_H88\_*} and \texttt{Ergebnis\_H89\_*}.}
\label{tab:repro-act}
\end{table}

\paragraph{Tools.} Numerical work uses Python (NumPy, SciPy, sympy for symbolic steps). Wolfram Language is used for exact
and symbolic derivations and as an independent second system, as in the derivation of the field equations of W1 and the
recomputation of the Big Bang admixture (\rndold{r86}). A separate Wolfram audit archive with input scripts and recorded
outputs for the algebraic core, the Dixon carrier, the Koide geometry and the OP5 comparison model is in
\nolinkurl{paper/preprint-de/audit/wolfram-v0.1.0/}. Round labels, cross references, guards and the ledger are checked with
\nolinkurl{tools/qs.py}.

\paragraph{Rebuilding.} The article is self-contained in \nolinkurl{paper/physics-standalone/}: \texttt{article.tex} with
\nolinkurl{sections/}, the TikZ figures in \nolinkurl{figures/} and the bibliography in \nolinkurl{refs/}. Run \texttt{pdflatex
article.tex} twice in that directory; scripts that generate figure data are in \nolinkurl{scripts/}. The Companion is built
from \nolinkurl{sources/research/v480-companion.tex} with that directory as working directory and several pdfLaTeX runs, as
described in \nolinkurl{sources/research/README.md}.
